% Posiciones cualitativas propuestas; PI-01 a PI-04 son roles de la misma persona.
\begin{tikzpicture}[x=0.76cm,y=0.68cm,
parte/.style={rectangle,rounded corners=2pt,fill=uteqverde,text=white,inner sep=2pt,font=\scriptsize\bfseries}]
\fill[gray!6] (0,0) rectangle (10,10);
\fill[green!9] (5,5) rectangle (10,10);
\draw[thick] (0,0) rectangle (10,10);
\draw[dashed] (5,0)--(5,10);\draw[dashed] (0,5)--(10,5);
\node[font=\small\itshape,align=center] at (2.5,9.4) {Mantener satisfecho};
\node[font=\small\itshape,align=center] at (7.5,9.4) {Gestionar de cerca};
\node[font=\small\itshape] at (2.5,0.7) {Monitorear};
\node[font=\small\itshape] at (7.5,0.7) {Mantener informado};
\draw[-{Stealth}] (0,0)--(10.5,0) node[right] {Interés};
\draw[-{Stealth}] (0,0)--(0,10.5) node[above] {Poder};
\node[below,font=\scriptsize] at (0,0) {Bajo};
\node[below,font=\scriptsize] at (10,0) {Alto};
\node[left,font=\scriptsize] at (0,10) {Alto};
\node[parte] at (9.1,8.3) {PI-01};
\node[parte] at (6.1,7.1) {PI-02};
\node[parte] at (7.6,7.1) {PI-03};
\node[parte] at (9.1,7.1) {PI-04};
\node[parte] at (8.5,2.4) {PI-05};
\node[parte] at (5.8,2.4) {PI-06};
\node[parte] at (6.1,8.3) {PI-07};
\node[parte] at (7.6,8.3) {PI-08};
\end{tikzpicture}
